\documentclass[10pt,a4paper]{amsart}
\usepackage[margin=19mm]{geometry}
\usepackage{amsmath,amssymb,mathtools}
\usepackage[scr=rsfs,cal=euler]{mathalfa}
\usepackage{xfrac}
\usepackage[unicode,hidelinks,bookmarks=false]{hyperref}
\newcommand{\ie}{i.e.\ }
\newcommand{\cf}{cf.\ }
\newcommand{\resp}{respectively\ }
\newcommand{\Subsection}{\subsection}
\providecommand{\nobreakdash}{\nobreak\hbox{-}}
\setlength{\emergencystretch}{2em}
\multlinegap0pt
\usepackage{bm}
\usepackage{bbm}
\usepackage{dsfont}
\usepackage{xspace}
\usepackage{cancel}
\usepackage{mathtools}
\usepackage{amsthm}
\makeatletter
\@ifundefined{theorem}{\newtheorem{theorem}{Theorem}}{}
\@ifundefined{proposition}{\newtheorem{proposition}[theorem]{Proposition}}{}
\makeatother
\newcommand{\End}{\mathrm{End}}
\newcommand{\E}{\mathcal{E}}
\newcommand{\Hom}{\mathrm{Hom}}
\newcommand{\Id}{\mathrm{Id}}
\newcommand{\dbar}{\overline{\partial}}
\newcommand{\e}{\boldsymbol{\mathrm{e}}}
\newcommand{\mU}{\mathcal{U}}
\newcommand{\norme}[1]{\left\| #1 \right\|}
\newcommand{\scal}[2]{\left< #1,#2 \right>}
\newcommand{\tend}[2]{\underset{#1 \rightarrow #2}{\longrightarrow}}
\renewcommand{\geq}{\geqslant}
\renewcommand{\i}{\boldsymbol{\mathrm{i}}}
\renewcommand{\leq}{\leqslant}
\title[Perturbations of Vector Bundle whose Curvature Form Solves a]{Perturbations of Vector Bundle whose Curvature Form Solves a Polynomial Equation}
\author{R\'emi Delloque}
\date{Source: arXiv:2507.01534v1 (CC BY 4.0)}
\begin{document}
\maketitle

\noindent\emph{Source and licence.} Perturbations of Vector Bundle whose Curvature Form Solves a Polynomial Equation. Version arXiv:2507.01534v1 (CC BY 4.0), distributed under the Creative Commons Attribution 4.0 International licence, as recorded in the arXiv licence metadata for this exact version and shown on the abstract page https://arxiv.org/abs/2507.01534v1. Full rights notice: licenses/arxiv-2507.01534-CC-BY-4.0.txt.\par\smallskip

\noindent\textit{Editorial scope.} Counted unit: the Proposition whose source label is \texttt{PRO:Convergence exp(txi)b} in the article (source0.tex lines 1222--1231, source label \texttt{PRO:Convergence exp(txi)b}), delivered together with the article's own complete original proof of it (source0.tex lines 1232--1385, one \texttt{proof} environment of 154 lines). No mathematics is written, repaired or re-derived: statement and proof are the article's own text line for line. Required context retained uncounted: none. Every definition, display and label of those lines is retained unchanged. Omitted from this package and not redistributed: 1--1221, 1386--3152. Nothing inside a retained excerpt is omitted or altered: every retained body is delivered exactly as the article's lines are written, so no omission receipt of any kind accompanies this package. Citations: the retained excerpts carry no citation command at all, so no bibliography entry of the article is transcribed and no reference is redistributed. Reference adaptation: none was needed and none was made; every reference inside the delivered unit resolves inside it, so no unresolved reference or citation is produced. Printed slips kept literal: no printed word, symbol, hypothesis or number of the counted statement or of its proof was altered. Layout and class: the article's own class is \texttt{article} with packages including geometry, chngpage, babel, amsmath, amsthm, amssymb, tikz-cd, dsfont, picture, float, graphicx, listings, hyperref, verbatim; the wrapper is the lane's pinned \texttt{amsart} editorial wrapper at 10pt on A4, which restates the theorem-like environments the retained text needs and adds the article's own macro definitions that the retained text needs (\texttt{\textbackslash E}, \texttt{\textbackslash End}, \texttt{\textbackslash Hom}, \texttt{\textbackslash Id}, \texttt{\textbackslash dbar}, \texttt{\textbackslash e}, \texttt{\textbackslash geq}, \texttt{\textbackslash i}, \texttt{\textbackslash leq}, \texttt{\textbackslash mU}, \texttt{\textbackslash norme}, \texttt{\textbackslash scal}, \texttt{\textbackslash tend}), with extra wrapper packages (bm, bbm, dsfont, xspace, cancel, mathtools). The delivered numbering is the unit's own. Assets: no retained span contains a figure environment, a graphics-inclusion command, a PDF-page inclusion, a table or a TikZ picture; no image file is copied into this package, the package declares no asset dependency and no image file is redistributed with it. Licence: version arXiv:2507.01534v1 of this article is distributed under the Creative Commons Attribution 4.0 International licence, as recorded in the arXiv licence metadata for this exact version and shown on the abstract page https://arxiv.org/abs/2507.01534v1. A full rights notice is delivered at licenses/arxiv-2507.01534-CC-BY-4.0.txt. Characters: every retained excerpt is pure ASCII, so the delivered engine, XeLaTeX with its default Latin Modern text font, reports no missing character.
\begin{proposition}\label{PRO:Convergence exp(txi)b}
    Up to shrinking $B$ and $\mU_d$, for all $\zeta \in \mU_d$, all $b \in B$ and all $\xi = \sum_{\lambda \in \mathrm{Sp}(\xi)} \lambda\Pi_\lambda$ in $\i\frak{k}$, we have equivalence between,
    \begin{enumerate}
        \item For all $\lambda$, $(F_{\xi,\lambda},\dbar_b) \subset \E_b$ is holomorphic (in other words, they define a filtration of $\E_b$ by admissible sub-bundles),
        \item For all $t \geq 0$, $\e^{t\xi} \cdot b \in B_{4,0}$ and it converges in $B_{4,0}$ when $t \rightarrow +\infty$,
        \item There is are sequences $t_p \rightarrow +\infty$, $b_p \rightarrow b$ and $\xi_p \rightarrow \xi$ of elements of $\i\frak{k}$ such that $(\e^{t_p\xi_p} \cdot b_p)$ is bounded in $V$.
        \item $\norme{\e^{t\xi} \cdot \dbar_b}_{L^2}$ is bounded when $t \rightarrow +\infty$.
    \end{enumerate}
    Moreover, in this case, if we set $b_\infty = \lim_{t \rightarrow +\infty} \e^{t\xi} \cdot b \in B_{4,0}$, all the embeddings $(G_{\xi,\lambda},\dbar_{b_\infty}) \subset \E_{b_\infty}$ are holomorphic and for all $t$, $\norme{\e^{t\xi} \cdot b} \leq \norme{b}$. In particular, $b_\infty \in B$.
\end{proposition}

\begin{proof}\ \\
\noindent\framebox{$1 \Rightarrow 2$} Assume the first point. Let $\mathrm{Sp}(\xi) = \{\lambda_1,\ldots,\lambda_m\}$ with $\lambda_1 < \lambda_2 < \ldots < \lambda_m$. It means that in the orthogonal decomposition
$$
E = \bigoplus_{k = 1}^m G_{\xi,\lambda_k},
$$
we have
$$
\Phi(b) = \dbar_b - \dbar_0 =
\begin{pmatrix}
    \gamma_{11} & \gamma_{12} & \hdots & \gamma_{1m}\\
    0 & \gamma_{22} & \hdots & \gamma_{2m}\\
    \vdots & \vdots & \ddots & \vdots\\
    0 & 0 & \hdots & \gamma_{mm}
\end{pmatrix}
$$
because each $F_{\xi,\lambda_k} = \bigoplus_{i = 1}^k G_{\xi,\lambda_i} \subset E$ is preserved by both $\dbar_0$ and $\dbar_b$. Let $\Id_k$ be the identity of $G_{\xi,\lambda_k}$. We have, for all $t$,
\begin{align*}
    &\e^{t\xi} \cdot \Phi(b)\\
    =\ & \e^{t\xi}\Phi(b)\e^{-t\xi}\\
    =\ &
    \begin{pmatrix}
        \e^{\lambda_1t}\Id_1 & 0 & \hdots & 0\\
        0 & \e^{\lambda_2t}\Id_2 & \hdots & 0\\
        \vdots & \vdots & \ddots & \vdots\\
        0 & 0 & \hdots & \e^{\lambda_mt}\Id_m
    \end{pmatrix}
    \begin{pmatrix}
        \gamma_{11} & \gamma_{12} & \hdots & \gamma_{1m}\\
        0 & \gamma_{22} & \hdots & \gamma_{2m}\\
        \vdots & \vdots & \ddots & \vdots\\
        0 & 0 & \hdots & \gamma_{mm}
    \end{pmatrix}
    \begin{pmatrix}
        \e^{-\lambda_1t}\Id_1 & 0 & \hdots & 0\\
        0 & \e^{-\lambda_2t}\Id_2 & \hdots & 0\\
        \vdots & \vdots & \ddots & \vdots\\
        0 & 0 & \hdots & \e^{-\lambda_mt}\Id_m
    \end{pmatrix}\\
    =\ &
    \begin{pmatrix}
        \gamma_{11} & \e^{-(\lambda_2 - \lambda_1)t}\gamma_{12} & \hdots & \e^{-(\lambda_m - \lambda_1)t}\gamma_{1m}\\
        0 & \gamma_{22} & \hdots & \e^{-(\lambda_m - \lambda_2)t}\gamma_{2m}\\
        \vdots & \vdots & \ddots & \vdots\\
        0 & 0 & \hdots & \gamma_{mm}
    \end{pmatrix}\\
    \tend{t}{+\infty}\ &
    \begin{pmatrix}
        \gamma_{11} & 0 & \hdots & 0\\
        0 & \gamma_{22} & \hdots & 0\\
        \vdots & \vdots & \ddots & \vdots\\
        0 & 0 & \hdots & \gamma_{mm}
    \end{pmatrix}.
\end{align*}
The decomposition $E = \bigoplus_{k = 1}^m G_{\xi,\lambda_k}$ is orthogonal and $\dbar_0$-holomorphic thus the associated decomposition
$$
\Omega^{0,1}(X,\End(E)) = \bigoplus_{i,j = 1}^m \Omega^{0,1}(X,\Hom(G_{\xi,\lambda_j},G_{\xi,\lambda_i}))
$$
is orthogonal for the Hermitian product $\scal{\cdot}{\cdot}_{L^2}$. We deduce that $t \mapsto \norme{\e^{t\xi} \cdot \Phi(b)}_{L^2}$ is non-increasing. Let $B'$ be an open ball in $V$ such that $B \subsetneq B' \subsetneq B_{4,0}$. As $\Phi$ is an embedding and $\partial B'$ is compact and does not contain $0$,
$$
a = \inf_{\partial B'} \norme{\Phi(b)}_{L^2}
$$
is positive. Up to shrinking $B$, we may assume that for all $b \in B$, $\norme{\Phi(b)}_{L^2} < a$. In particular, if $b \in B$, for all $t \geq 0$, $\norme{\e^{t\xi} \cdot \Phi(b)}_{L^2} < a$. It implies that whenever $\e^{t\xi} \cdot b \in B_{4,0}$, it is not in $\partial B'$ (recall that $\Phi$ is $G$-invariant). Thus, it stays in the compact subset $\overline{B'}$ of $B_{4,0}$. In particular, it implies that for all non-negative time, $\e^{t\xi} \cdot b \in \overline{B'} \subset B_{4,0}$ and it converges because $\Phi(\e^{t\xi} \cdot b)$ converges and $\Phi_{|\overline{B'}}$ is injective and proper.

Moreover, the limit structure $\dbar_{b_\infty}$ is diagonal in the decomposition $E = \bigoplus_{k = 1}^m G_{\xi,\lambda_k}$, each $G_{\xi,\lambda_k}$ is preserved hence the holomorphic decomposition
$$
\E_{b_\infty} = \bigoplus_{k = 1}^m (G_{\xi,\lambda_k},\dbar_{b_\infty}).
$$\\

\noindent\framebox{$2 \Rightarrow 3$ and $4$} Trivial.\\

\noindent\framebox{$3 \Rightarrow 2$} First of all, the set of all $\eta$ which are diagonal in the decomposition
$$
\E_0 = \bigoplus_{k = 1}^m (G_{\xi,\lambda_k},\dbar_0)
$$
is the Lie algebra $\frak{t}$ of a complex torus $T$ in $G$. Therefore, it is contained in the Lie algebra $\frak{t}_{\max}$ of a maximal torus $T_{\max}$ of $G$. It can be obtained thanks to an orthogonal and $\dbar_0$-holomorphic decomposition of each $(G_{\xi,\lambda_k},\dbar_0)$ for example. By assumption, $\xi \in \frak{t} \subset \frak{t}_{\max}$. Since $\xi_p \rightarrow \xi$, we can diagonalise it as
$$
\xi_p = u_p\xi_p'u_p^\dagger,
$$
where $u_p \in K$ and $\xi_p' \in \frak{t}_{\max}$. Moreover, as $\xi_p \rightarrow \xi \in \frak{t}_{\max}$, $\xi_p' \rightarrow \xi$ and we may assume that $u_p \rightarrow \Id_E$. Let us write, in the orthogonal decomposition $E = \bigoplus_{k = 1}^m G_{\xi,\lambda_k}$,
$$
\xi_p' =
\begin{pmatrix}
    \lambda_{p,1} & 0 & \hdots & 0\\
    0 & \lambda_{p,2} & \hdots & 0\\
    \vdots & \vdots & \ddots & \vdots\\
    0 & 0 & \hdots & \lambda_{p,m}
\end{pmatrix}.
$$
There is no reason to believe that $\xi_p'$ is in $\frak{t}$ so each $\lambda_{p,k}$ is a holomorphic Hermitian endomorphism of $(G_{\xi,\lambda_k},\dbar_0)$ but it needs not be a homothety. In other words, it is block diagonal but not diagonal. As $\xi_p' \rightarrow \xi$, we have $\lambda_{p,k} \rightarrow \lambda_k\Id_k$ for all $k$. Let for all $p$, $b_p' = u_p^\dagger \cdot b_p$ and let us set
$$
b =
\begin{pmatrix}
    b_{11} & b_{12} & \hdots & b_{1m}\\
    b_{21} & b_{22} & \hdots & b_{2m}\\
    \vdots & \vdots & \ddots & \vdots\\
    b_{m1} & b_{m2} & \hdots & b_{mm}
\end{pmatrix}, \qquad
b_p' = 
\begin{pmatrix}
    b_{p,11} & b_{p,12} & \hdots & b_{p,1m}\\
    b_{p,21} & b_{p,22} & \hdots & b_{p,2m}\\
    \vdots & \vdots & \ddots & \vdots\\
    b_{p,m1} & b_{p,m2} & \hdots & b_{p,mm}
\end{pmatrix}.
$$
As $u_p \rightarrow \Id_E$, for all $i,j$, $b_{p,ij} \rightarrow b_{ij}$. We have, for all $p$,
$$
\e^{t_p\xi_p} \cdot b_p = u_p\e^{t_p\xi_p'} \cdot b_p' = u_p \cdot
\begin{pmatrix}
    \e^{t_p\lambda_{p,1}}b_{p,11}\e^{-t_p\lambda_{p,1}} & \e^{t_p\lambda_{p,1}}b_{p,12}\e^{-t_p\lambda_{p,2}} & \hdots & \e^{t_p\lambda_{p,1}}b_{p,1m}\e^{-t_p\lambda_{p,m}}\\
    \e^{t_p\lambda_{p,2}}b_{p,21}\e^{-t_p\lambda_{p,1}} & \e^{t_p\lambda_{p,2}}b_{p,22}\e^{-t_p\lambda_{p,2}} & \hdots & \e^{t_p\lambda_{p,2}}b_{p,2m}\e^{-t_p\lambda_{p,m}}\\
    \vdots & \vdots & \ddots & \vdots\\
    \e^{t_p\lambda_{p,m}}b_{p,m1}\e^{-t_p\lambda_{p,1}} & \e^{t_p\lambda_{p,m}}b_{p,m2}\e^{-t_p\lambda_{p,2}} & \hdots & \e^{t_p\lambda_{p,m}}b_{p,mm}\e^{-t_p\lambda_{p,m}}
\end{pmatrix}.
$$
By assumption, this sequence is bounded so for all $i,j$, $(\e^{t_p\lambda_{p,i}}b_{p,ij}\e^{-t_p\lambda_{p,j}})$ is bounded. Let $i > j$, so $\lambda_i > \lambda_j$. Set $\varepsilon_{p,i} = \lambda_{p,i} - \lambda_i\Id_i \rightarrow 0$ and $\varepsilon_{p,j} = \lambda_{p,j} - \lambda_j\Id_j \rightarrow 0$. We have
\begin{align*}
    \norme{\e^{t_p\lambda_{p,i}}b_{p,ij}\e^{-t_p\lambda_{p,j}}}_{L^2} & = \norme{\e^{t_p\varepsilon_{p,i}}\e^{(\lambda_i - \lambda_j)t_p}b_{p,ij}\e^{-t_p\varepsilon_{p,j}}}_{L^2}\\
    & \geq \norme{\e^{-t_p\varepsilon_{p,i}}}^{-1}\norme{\e^{t_p\varepsilon_{p,j}}}^{-1}\norme{\e^{(\lambda_i - \lambda_j)t_p}b_{p,ij}}_{L^2}\\
    & \geq \e^{-t_p\norme{\varepsilon_{p,i}}}\e^{-t_p\norme{\varepsilon_{p,j}}}\e^{(\lambda_i - \lambda_j)t_p}\norme{b_{p,ij}}_{L^2}\\
    & \geq \e^{\frac{1}{2}(\lambda_i - \lambda_j)t_p}\norme{b_{p,ij}}_{L^2} \textrm{ for $p$ large enough.}
\end{align*}
Since this sequence is bounded, we must have $\norme{b_{p,ij}}_{L^2} \rightarrow 0$ so $b_{ij} = 0$. It holds for all $i > j$ so $b$ is upper diagonal. Moreover, the decomposition
$$
\Omega^{0,1}(X,\End(E)) = \bigoplus_{i,j = 1}^m \Omega^{0,1}(X,\Hom(G_j,G_i))
$$
is orthogonal. It implies that
$$
t \mapsto \norme{\e^{t\xi} \cdot b}_{L^2} = \sum_{i \leq j} \e^{t(\lambda_i - \lambda_j)}\norme{b_{ij}}_{L^2}
$$
decreases and $b_\infty = \lim_{t \rightarrow +\infty} \e^{t\xi} \cdot b$ exists and is given by the diagonal part of $b$. In particular, for all $t$, $\norme{\e^{t\xi} \cdot b}_{L^2} \leq \norme{b}_{L^2}$ and $\norme{b_\infty}_{L^2} \leq \norme{b}_{L^2}$ so $\e^{t\xi} \cdot b$ and $b_\infty$ lie in $B \subset B_{4,0}$.\\

\noindent\framebox{$4 \Rightarrow 1$} Write, in the orthogonal decomposition $E = \bigoplus_{k = 1}^m G_{\xi,\lambda_k}$,
$$
\Phi(b) =
\begin{pmatrix}
\gamma_{11} & \gamma_{12} & \hdots & \gamma_{1m}\\
\gamma_{21} & \gamma_{22} & \hdots & \gamma_{2m}\\
\vdots & \vdots & \ddots & \vdots\\
\gamma_{m1} & \gamma_{m2} & \hdots & \gamma_{mm}
\end{pmatrix}.
$$
As before, we have, for all $t$,
$$
\e^{t\xi} \cdot \Phi(b) =
\begin{pmatrix}
\gamma_{11} & \e^{-(\lambda_2 - \lambda_1)t}\gamma_{12} & \hdots & \e^{-(\lambda_m - \lambda_1)t}\gamma_{1m}\\
\e^{(\lambda_2 - \lambda_1)t}\gamma_{21} & \gamma_{22} & \hdots & \e^{-(\lambda_m - \lambda_2)t}\gamma_{2m}\\
\vdots & \vdots & \ddots & \vdots\\
\e^{(\lambda_m - \lambda_1)t}\gamma_{m1} & \e^{(\lambda_m - \lambda_2)t}\gamma_{m2} & \hdots & \gamma_{mm}
\end{pmatrix}.
$$
We easily see that if $\norme{\e^{t\xi} \cdot \dbar_b}_{L^2}$ is bounded for large $t$, then $\gamma_{ij} = 0$ when $i > j$. We deduce that each $F_{\xi,\lambda_k}$ is preserved by $\dbar_b = \dbar_0 + \Phi(b)$ thus $(F_{\xi,\lambda_k},\dbar_b) \subset \E_b$ is holomorphic.
\end{proof}

\end{document}
